\documentclass[preprint,12pt,authoryear]{elsarticle}
\usepackage{fontspec}
\setmainfont{Liberation Serif}
\setsansfont{Liberation Sans}
\setmonofont{Liberation Mono}
\usepackage{polyglossia}
\setdefaultlanguage{russian}
\setotherlanguage{english}
\usepackage{amsmath,amssymb}
\usepackage{booktabs}
\usepackage{graphicx}
\usepackage{microtype}
\usepackage{multirow}
\usepackage{longtable}
\usepackage{array}
\usepackage[section]{placeins}
\usepackage{pgfplots}
\pgfplotsset{compat=1.17}
\usepackage[hidelinks]{hyperref}
\usepackage{lineno}
\biboptions{round,authoryear}
\setlength{\emergencystretch}{3em}

% \journal{Decision Support Systems}  % restore when actually submitting
\makeatletter
\def\ps@pprintTitle{\let\@oddhead\@empty\let\@evenhead\@empty
  \def\@oddfoot{\footnotesize\itshape Рукопись в подготовке\hfill\today}\let\@evenfoot\@oddfoot}
\makeatother

\newcommand{\dE}{\Delta E}
\newcommand{\pp}{\,п.\,п.}

\abstracttitle{Аннотация}
\keywordtitle{Ключевые слова}
\begin{document}

\begin{frontmatter}

\title{Операционная автономность систем поддержки принятия решений на основе LLM при отключении поставщиков: общие ошибки, ограничения «одна модель от разработчика» и семейства моделей}

\author{А.~В. Стрижевский (strizhevskiyav@student.bmstu.ru)}
\author{Е.~Д. Верховых (verkhovyykhed@student.bmstu.ru)}
\author{Д.~М. Парамонов (paramonovdm@student.bmstu.ru)}
\author{М.~И. Колосов (kolosovm@bmstu.ru)}
\address{МГТУ им. Н.~Э. Баумана, Москва, Россия}

\begin{abstract}
Системы поддержки принятия решений могут объединять несколько больших языковых моделей (LLM), а доступ к ним может быть утрачен группами: ко всем моделям одного разработчика или ко всем моделям, доступным только через API. Мы предлагаем стресс-тест отключения, который рассматривает портфель LLM как систему компонентов и измеряет его операционную автономность мерами значимости из теории надёжности: приростом ошибки при групповом отключении, его относительной формой (risk achievement worth) и совместной значимостью. Гипотезы были зафиксированы во внутреннем протоколе и проверены на синтетических панелях до анализа реальных ответов. По ответам 70 моделей на 14\,038 вопросов MMLU, репликации на MMLU-Pro, 13 локальным запускам и парным данным о 349 моделях рейтинга получено: для основного портфеля на MMLU потеря любого одного разработчика стоит не более 0.6 п.~п. ошибки, тогда как ограничение системы открытыми моделями стоит в этой панели 2.3--2.9 п.~п. на MMLU и около 5.2 п.~п. на MMLU-Pro; для открытых моделей удаление означает вынужденный отказ, а не утрату доступа. Модели одного разработчика совпадают в неверных ответах сверх сложности вопросов: устойчиво на MMLU и слабо на MMLU-Pro. Ни один из трёх проверенных методов отбора не превзошёл ограничение «одна модель от разработчика», а оценка в предположении независимости завышает стоимость отключения примерно в пять раз. Дообученные модели сохраняют значительную часть структуры ошибок базовой модели. При 8--12 млрд параметров отставания российских моделей по точности на MMLU не обнаружено, что не доказывает их равенства. Метод измеряет устойчивость точности ответов, а не операционный риск, который требует проверки на отраслевых задачах. Результаты подкрепляют тестирование отключений и ограничение по разработчикам; подсчёт семейств моделей, а не брендов, остаётся гипотезой для проектирования, которую прямая проверка не подтвердила и не опровергла.
\end{abstract}

\begin{keyword}
системы поддержки принятия решений \sep большие языковые модели \sep ансамбли моделей \sep коррелированные ошибки \sep меры значимости в теории надёжности \sep технологический суверенитет
\end{keyword}

\end{frontmatter}


\section{Введение}

Системы поддержки принятия решений (СППР) используют большие языковые модели (LLM) в качестве компонентов, которые отвечают на вопросы, классифицируют документы и готовят проекты рекомендаций. Надёжная архитектура редко опирается на одну модель: ответы нескольких моделей агрегируются голосованием или обучаемым комбинатором, либо запросы маршрутизируются между моделями \citep{jiang2023blender,wang2025moa,chen2024frugalgpt,ong2025routellm}. У такой архитектуры есть зависимость, которую оценки средней точности не учитывают. Модели поступают от небольшого числа разработчиков, и доступ к ним может быть утрачен группами. Разработчик может вывести из эксплуатации семейство моделей, изменить условия использования или прекратить обслуживание какой-либо юрисдикции; все модели, доступные только через API, могут стать недоступны одновременно, что является формой зависимости от поставщика (vendor lock-in), давно обсуждаемой применительно к облачным сервисам \citep{oparamartins2016}. Поэтому для организаций критически важных отраслей выбор моделей является также вопросом технологического суверенитета, который литература по инновационной политике определяет как способность государства обеспечивать технологии, которые оно считает критически важными, разрабатывая их или получая из других экономических зон без односторонней структурной зависимости \citep{edler2023}. Для части операторов это и вопрос соответствия требованиям: в России компании энергетики и топливно-энергетического комплекса относятся к субъектам критической информационной инфраструктуры \citep{fz187}, а органам власти и ряду заказчиков с 1 января 2025 г. запрещено использовать иностранное программное обеспечение на принадлежащих им значимых объектах этой инфраструктуры \citep{ukaz166}.

Мы изучаем надёжность таких систем на уровне компонентов: каждая LLM~--- компонент, ответы которого могут быть ошибочными, а надёжность системы~--- точность агрегированного ответа. В статье рассматриваются три вопроса, с которыми сталкивается проектировщик СППР. (i) Какую часть точности сохраняет агрегированная система при отключении группы моделей? (ii) Как следует выбирать модели и правило агрегации, чтобы ограничить эти потери? (iii) Что даёт новая, например отечественная, модель и ведёт ли себя дообученная модель на основе зарубежной базовой модели как независимый компонент? Мы отвечаем на них с помощью \emph{стресс-теста отключения}, который рассматривает портфель LLM как систему компонентов и использует меры значимости, разработанные в теории надёжности для компонентов, отказывающих совместно \citep{birnbaum1969,vesely1983,borgonovo2001,honglie1993}. В терминах планирования энергосистем это проверка портфеля по критерию N$-$1 (отключается любой один разработчик), расширенная на пары разработчиков (N$-$2) и на отказ по общей причине~--- отключение всех моделей, доступных только через API.

Исследование определяют три особенности. Во-первых, нулевая модель учитывает сложность вопроса и популярность каждого неверного варианта ответа, так что общие ошибки не смешиваются с общими трудными вопросами \citep{eckhardt1985}. Во-вторых, гипотезы, решающие правила и параметры анализа были зафиксированы в письменном внутреннем протоколе с записанным хешем и проверены на синтетических панелях с известной истиной до анализа какого-либо реального ответа; последующие изменения и анализы post hoc (апостериорный анализ) перечислены с указанием времени их проведения (\ref{app:amend}). В-третьих, анализ использует публичные записи оценок и опубликованные нами локальные запуски, поэтому результаты могут быть воспроизведены без доступа к моделям, размещённым у поставщиков.

Наш вклад состоит в следующем:
\begin{enumerate}
\item Стресс-тест отключения для портфелей LLM, который сопоставляет потери при групповом отключении, их относительную форму и совместную значимость с мерами значимости вероятностного анализа безопасности, с синтетической проверкой тестов и покрытия их интервалов.
\item Сопоставимая (like-for-like) стоимость ограничения системы открытыми моделями на двух бенчмарках и сравнение с прогнозами стоимости отключения, основанными на независимости и на сложности вопросов.
\item Свидетельства, воспроизведённые (слабее) на втором бенчмарке, где интервал по методу складного ножа по разработчикам включает ноль, того, что общие ошибки на уровне разработчика сохраняются в тесте, калиброванном по сложности, при нескольких мерах согласия после выравнивания точности, и того, что ни один из трёх проверенных методов отбора не превзошёл ограничение «одна модель от разработчика»: на MMLU его ошибка в наихудшем случае на отложенных данных (11.65\%) равна ошибке жадного отбора по точности (11.65\%) и немного ниже, чем у отбора, ориентированного на устойчивость (11.81\%); выигрыш ограничения появляется только на MMLU-Pro (11.10\% против 13.10\%; приближённые интервалы).
\item Количественная оценка наследования ошибок от базовых моделей к дообученным~--- прежде всего на уровне вопросов с нулевой моделью сложности для шести связанных локальных пар; её подкрепляет анализ пар открытого рейтинга, в котором прямая связь дообучения добавляет общие ошибки сверх класса «архитектура и размер» семейства моделей.
\end{enumerate}

\section{Обзор литературы}

\paragraph{Комбинирование LLM} Выходы LLM комбинируются ранжированием и слиянием \citep{jiang2023blender}, многоуровневыми архитектурами, в которых модели уточняют ответы друг друга \citep{wang2025moa}, каскадами, которые сначала обращаются к более дешёвым моделям \citep{chen2024frugalgpt}, и обучаемыми маршрутизаторами \citep{ong2025routellm}. Эти методы оптимизируют качество и стоимость при доступности всех моделей. LLM также непосредственно входят в исследования СППР, например для преобразования объяснений моделей в повествования \citep{martens2025}. Устойчивость таких комбинаций к потере их компонентов, насколько нам известно, не изучалась.

\paragraph{Коррелированные ошибки LLM} \citet{kim2025correlated} измерили согласие в неверных ответах для 349 открытых моделей в Open LLM Leaderboard и 71 модели в HELM и обнаружили, что модели одной компании и одной архитектуры чаще совпадают, когда обе ошибаются. Их мера сравнивает согласие с постоянным случайным уровнем, и авторы указывают отсутствие контроля сложности в качестве ограничения. \citet{goel2025great} предложили меру согласия с поправкой на случайность, обусловленную точностью (CAPA), и показали, что похожие модели в роли судей отдают предпочтение друг другу; пары моделей одного разработчика они исключают как смешивающий фактор, а не изучают. Наш тест на уровне разработчика воспроизводит направление эффекта из \citet{kim2025correlated}; его дополнения~--- калибровка по смоделированному миру с общей сложностью и популярностью неверных вариантов, минимальный обнаружимый эффект, второй бенчмарк и сравнение мер согласия.

\paragraph{Происхождение (родословная) LLM} Происхождение оставляет следы в выходах моделей: филогении можно восстановить по сходству выходов \citep{yax2025phylolm}, а ErrorTrace относит модель к её базовому семейству по множеству её ошибок \citep{zang2025errortrace}. Эти работы используют происхождение для атрибуции; в них не сообщается, какую часть структуры отказов базовой модели сохраняет производная модель по сравнению с несвязанными моделями сопоставимой точности, и что это означает для ансамблей.

\paragraph{Разнообразие ансамблей и качество бенчмарков} Разнообразие~--- классический принцип построения ансамблей классификаторов \citep{kuncheva2003,brown2005}; его меры вычисляются на фиксированной команде и смешивают все источники совпадающих отказов. \citet{turkmen2026} жадно отбирают ансамбли LLM по взаимной информации при коррелированных ошибках в условиях полной доступности. Метки бенчмарков несовершенны: повторная разметка выявила ошибки примерно в 6.5\% вопросов MMLU \citep{gema2025mmlu}, что важно для мер согласия, поскольку сильные модели совпадают в «неверных» ответах на вопросах с ошибочной разметкой.

\paragraph{Теория надёжности} Меры значимости ранжируют компоненты по их влиянию на риск системы: значимость Бирнбаума \citep{birnbaum1969}, risk achievement worth (RAW; риск при отказавшем компоненте относительно номинального; \citealp{vesely1983}), мера дифференциальной значимости, которая, в отличие от risk achievement worth и значимости Фасселла--Весели, аддитивна по группам \citep{borgonovo2001}, и совместная значимость~--- взаимодействие второго порядка двух компонентов \citep{honglie1993}. Метрики живучести (resilience) разделяют потери после нарушения и последующее восстановление \citep{henry2012,hosseini2016}. Отказы по общей причине~--- это множественные отказы компонентов, вызванные общими первопричинами \citep{iaea1992ccf}. В многоверсионном программном обеспечении независимо разработанные версии отказывают независимо, только если вероятность отказа одинакова для всех входных данных; если она меняется от входа к входу, совместные отказы превышают предсказание модели независимости \citep{eckhardt1985}, что подтвердил эксперимент \citep{knight1986}. Разработка версий разными методологиями снижает вероятность одновременного отказа и теоретически может даже превзойти независимость, но не гарантирует заданного уровня надёжности на практике \citep{littlewood1989}. Мы используем этот понятийный аппарат напрямую: нулевая модель~--- базовый уровень совпадающих отказов из-за общей сложности входных данных, эффекты разработчика и семейства~--- зависимости по общей причине сверх него, а меры отключения~--- меры значимости, применённые к группам моделей.

\section{Постановка задачи и меры}

\subsection{Условия}
Панель LLM отвечает на $N$ вопросов с выбором ответа с правильными ключами $y_k$. Модель $i$ даёт ответ $a_{ik}$. \emph{Портфель} $P$~--- множество из $n$ моделей, объединённых правилом агрегации; его ошибка $E(P)$~--- доля вопросов, на которые агрегированный ответ неверен. \emph{Сценарий отключения} $W$ удаляет группу моделей: все модели одного разработчика или все модели, доступные только через API (модели, веса которых не опубликованы для локального развёртывания). При отключении система заново подгоняет своё правило на оставшихся моделях. Для моделей, доступных только через API, отключение означает утрату доступа. Для разработчиков открытых моделей это более слабый сценарий~--- утрата будущих версий и поддержки, поскольку опубликованные веса можно хранить локально. Тем не менее наш стресс-тест в обоих случаях удаляет модели: для моделей, доступных только через API, вычисленная потеря является следствием утраты доступа, а для открытых моделей это условный сценарий вынужденного отказа от модели (например, по лицензионным, регуляторным соображениям или соображениям безопасности), а не неизбежное следствие прекращения поддержки разработчиком.

\subsection{Меры отключения}
Следуя мерам значимости вероятностного анализа безопасности, мы приводим
\begin{align}
\dE_W &= E(P\setminus W) - E(P), \\
\mathrm{RAW}_W &= E(P\setminus W)/E(P), \\
\mathrm{JI}(W_1,W_2) &= E(P\setminus(W_1\cup W_2)) - E(P\setminus W_1) - E(P\setminus W_2) + E(P).
\end{align}
$\dE_W$~--- основная мера; она соответствует слагаемому потерь в метриках живучести \citep{henry2012}; восстановление путём замены не моделируется. В классификации \citet{hosseini2016} меры отключения являются, таким образом, статическими мерами поглощающей способности системы, основанными на показателях функционирования. $\dE_W$ и $\mathrm{RAW}_W$~--- интервальная и относительная формы risk achievement worth, т.~е. прироста риска при удалении элемента или допущении его отказа \citep{vesely1983}. $\mathrm{JI}$~--- конечно-разностный аналог (по ошибке) совместной значимости для надёжности \citep{honglie1993}: он положителен, когда совместная потеря двух групп обходится дороже, чем сумма двух потерь. Если считать отключение отказом, а ошибку~--- дополнением надёжности до единицы, JI равна совместной значимости Хонга и Ли со знаком минус, так что положительная JI означает, что две группы взаимно заменяют друг друга. Эта мера нужна потому, что меры типа RAW не аддитивны по группам \citep{borgonovo2001}. Поскольку $\dE$ данного портфеля зависит от того, как этот портфель был отобран, мы также приводим сопоставимую \emph{стоимость суверенитета} $C = E(P_{\text{open}}) - E(P_{\text{cap}})$~--- ошибку лучшего портфеля из пяти открытых моделей минус ошибку портфеля с ограничением «одна модель от разработчика» при полной доступности, с парными бутстрэп-интервалами, а также её вариант $C'$, в котором портфель из пяти открытых моделей заменён лучшим открытым вариантом~--- отдельной моделью или голосованием, в зависимости от того, что ошибается реже.

\subsection{Нулевые модели}
N1 (независимость): каждая модель отвечает правильно со своей вероятностью, а неверные ответы равномерно распределены по дистракторам. N2 (условная независимость при заданном вопросе): правильность подчиняется однопараметрической модели ответа на задание с угадыванием, $P(\text{correct}) = g + (1-g)\,\sigma(\theta_i - b_k)$, где $g$~--- единица, делённая на число вариантов, а неверные ответы выбираются из объединённого распределения неверных ответов на данный вопрос $w_k$. N2 сохраняет сложность вопроса и популярность каждого неверного варианта и устраняет все прочие зависимости; она была выбрана вместо двухпараметрической модели в ходе синтетической проверки по заранее сформулированному правилу.

\subsection{Гипотезы}
\textbf{H1 (общие ошибки на уровне разработчика).} Пусть $S_{ij}$~--- число вопросов, на которые модели $i$ и $j$ дают одинаковый неверный ответ, с ожиданием по N2 $E_{ij} = \sum_k (1-p_{ik})(1-p_{jk}) \sum_o w_{ko}^2$, и превышение $x_{ij} = (S_{ij}-E_{ij})/N$. Статистика $T$~--- среднее по разработчикам от среднего превышения внутри разработчика минус среднее превышение пар разных разработчиков, подобранных к каждой внутренней паре по обеим точностям (в пределах 0.03 с расширением до 0.12, пока не найдётся пять подходящих пар) и по точному тексту промпта. H1 подтверждается, если односторонний $p$-value параметрического бутстрэпа при N2 ниже 0.05 и $T$ превышает минимальный обнаружимый эффект, установленный в моделировании; опровергается, если 90\% интервал по методу складного ножа (jackknife) по разработчикам лежит ниже этого эффекта.

\textbf{H2 (наиболее точное правило не является наиболее устойчивым).} Для портфелей из $n \in \{3,5,7,9\}$ наиболее точных моделей правило с наименьшей ошибкой при полной доступности отличается от правила с наименьшей ошибкой в наихудшем случае по сценариям отключения, причём не менее чем на 1 п.~п. по обоим показателям, как минимум для двух размеров. Опровергается, если при каждом размере победитель при полной доступности является также победителем в наихудшем случае или отстаёт от него не более чем на 0.5 п.~п.

\textbf{H3 (устойчивый отбор лучше ограничения разнообразия).} Устойчивый к сценариям отбор минимизирует $(1-\lambda)E(P) + \lambda \max_W E(P\setminus W)$ при $\lambda = 0.5$; сценарий, опустошающий портфель, засчитывается как ошибка 0.75 (угадывание). При оценке с исключением одного разработчика (leave-one-developer-out) он должен снизить ошибку в наихудшем случае на отложенных данных по сравнению с ограничением не более чем одной моделью от разработчика не менее чем на 1 п.~п. Опровергается, если ограничение столь же хорошо или лучше с точностью до 0.5 п.~п.

Правила агрегации: лучшая отдельная модель, голосование относительным большинством (ничьи разрешаются по суммарной точности моделей, проголосовавших за каждый вариант), взвешенное голосование по логарифму шансов, Dawid--Skene \citep{dawidskene1979} и стекинг. Для H3 сравниваются жадный отбор по точности, ограничение «одна модель от разработчика» и подстановочная (plug-in) эвристика взаимной информации в духе оценщика из \citet{turkmen2026}, но не идентичная ему.

\section{Данные и процедура}

\subsection{Бенчмарки и панели}
Проект Holistic Evaluation of Language Models \citep{liang2023helm} публикует ответ каждой оценённой модели на каждый вопрос. Основные данные~--- запуски HELM на MMLU \citep{hendrycks2021mmlu}: 14,038 вопросов с четырьмя вариантами ответа по 57 предметам, промпты в режиме 5-shot, температура 0. После устранения дубликатов снимков и псевдонимов по правилу, составленному по названиям моделей до изучения результатов, панель содержит 70 моделей от 18 разработчиков, из них 42 открытые и 28 доступных только через API (табл.~\ref{tab:data}). Промпты различаются только необязательной начальной инструкцией о формате; точный текст промпта является переменной подбора в H1, а 30 моделей с модальным промптом на каждом вопросе образуют панель для анализа чувствительности. Для репликации используются запуски HELM на MMLU-Pro \citep{wang2024mmlupro}: 1,000 вопросов с числом вариантов до десяти и ответами с цепочкой рассуждений, сгенерированными при температуре 1. Ответы извлекались из текста рассуждений; восемь из 66 моделей, у которых извлечённая правильность совпадала с собственной оценкой HELM менее чем на 95\% вопросов, были исключены из основной панели по заранее сформулированному правилу, что оставило 58 моделей от 15 разработчиков, 31 из них доступна только через API.

Мы также локально запустили 12 открытых моделей на точных промптах HELM MMLU с жадным декодированием, не более чем восемью выходными токенами и квантованными весами (в основном 4 бита, Q4\_K\_M; GigaChat3.1~--- 6 бит, GigaChat-20B~--- 5 бит, как их выпустили разработчики): восемь российских моделей (YandexGPT-5-Lite-8B, T-lite-it-2.1, GigaChat3.1-10B-A1.8B, GigaChat-20B-A3B, Vikhr-Nemo-12B, saiga\_\allowbreak nemo\_\allowbreak 12b, saiga\_\allowbreak yandexgpt\_\allowbreak 8b, Cotype-Nano) и четыре зарубежные опорные модели (Qwen2.5-7B, Qwen3-8B, Llama-3.1-8B, Mistral-Nemo-12B). Две российские модели декларируют зарубежную базовую модель (T-lite-it-2.1: Qwen3-8B; Vikhr-Nemo: Mistral-Nemo), а две являются дообученными моделями на основе базовой модели, указанной в названии модели (saiga\_nemo: Mistral-Nemo; saiga\_yandexgpt: YandexGPT-5-Lite). Тринадцатая, более крупная российская модель, T-pro-it-2.1 (32 млрд параметров, дообучение Qwen3-32B), была позднее запущена на всех вопросах (post hoc). Для наследования в масштабе популяции мы используем публичную таблицу пар из \citet{kim2025correlated} со статистиками согласия на MMLU-Pro для всех пар 349 моделей Open LLM Leaderboard, объединённую с самодекларируемым полем базовой модели из рейтинга; для 60,031 пары имеются метаданные обеих моделей.

\begin{table}[!htbp]
\centering
\caption{Источники данных.}
\label{tab:data}
\small
\resizebox{\textwidth}{!}{%
\begin{tabular}{@{}lrrrl@{}}
\toprule
Источник & Моделей & Вопросов & Разработчиков & Ответы \\
\midrule
HELM MMLU (основной) & 70 & 14,038 & 18 & на уровне вопросов, температура 0 \\
HELM MMLU-Pro (репликация) & 58 & 1,000 & 15 & на уровне вопросов, цепочка рассуждений, температура 1 \\
Локальные запуски (данное исследование) & 12 & 14,038 & 9 & на уровне вопросов, жадное декодирование, квантование \\
Open LLM Leaderboard \citep{kim2025correlated} & 349 & -- & -- & только парные статистики MMLU-Pro \\
\bottomrule
\end{tabular}}
\end{table}

\subsection{Оценивание}
Вопросы разбиваются на два блока перекрёстной подгонки (cross-fitting) (по предметам на MMLU, случайно на MMLU-Pro); отбор и подгонка выполняются на одном блоке, оценка~--- на другом, и оценки объединяют оба направления. Интервалы получены кластерным бутстрэпом по предметам на MMLU и по вопросам на MMLU-Pro (200 повторных выборок), причём каждый отбор и каждая подгонка повторяются внутри каждой повторной выборки; для совместной значимости используется бутстрэп-интервал, центрированный на оценке и расширенный в $\sqrt{G/(G-1)}$ раз. Для H1 используются 200 панелей параметрического бутстрэпа по N2 (100 в панелях для анализа чувствительности) и метод складного ножа (jackknife) по разработчикам. Основной портфель содержит $n=5$ моделей, отобранных жадно по точности при голосовании относительным большинством. Код анализа написан на Python (с помощью ИИ, см. раздел «Благодарности»); его корректность проверялась на синтетических панелях с известной истиной и путём воспроизведения каждой приведённой статистики по сохранённым выходным данным.

\subsection{Синтетическая проверка до анализа реальных данных}
До анализа какого-либо реального ответа тесты были выполнены на 600 синтетических панелях, сгенерированных с факторами разработчика и происхождения, сложностью вопросов, предпочтениями неверных вариантов и 3\% ошибочных ключей, на предварительной структуре из 64 моделей, близкой к реальной панели. Первая версия H1, основанная на потерях при отключении отдельных разработчиков, имела корректный фактический уровень значимости (размер теста), но мощность не более 0.12, поскольку удаление нескольких моделей из голосования 64 моделей почти не меняет его. Поэтому H1 была заменена описанной выше попарной статистикой. Дальнейшая проверка показала, что сравнение без подбора пар смещено при нулевой гипотезе, поскольку модели одного разработчика имеют сходные способности, а модели со сходными способностями совпадают в ответах на вопросах с ошибочной разметкой; это привело к выравниванию по точности. Окончательная H1 имеет уровень значимости 0.050, мощность 1.00 начиная со стандартного отклонения фактора 0.2 и минимальный обнаружимый эффект $T=0.0028$. Номинальные 95\% интервалы имели подтверждённое покрытие 0.91 для $\dE$ и 0.94 для JI; точечные оценки JI показали небольшое смещение вверх (16 из 22 промахов~--- ниже интервала). Диаграмма допуска (admission chart) восстанавливала истинные выигрыши с точностью до 0.1 п.~п. Для MMLU-Pro первая проверка уровня значимости не была пройдена (0.070); повторная проверка, принятая после этой неудачи (новые начальные значения генератора для уровня значимости, фиксированные сдвиги по предметам для покрытия), прошла проверку уровня значимости (0.050), но покрытие для $\dE$ составило 0.896, поэтому интервалы на MMLU-Pro были приближёнными, а решения там опираются на тест H1 и точечные оценки. Поскольку процентильные интервалы давали недостаточное покрытие, все приведённые ниже интервалы для $\dE$ откалиброваны на этих синтетических панелях (post hoc): обе полуширины относительно точечной оценки умножаются на коэффициент $k$, дающий покрытие 95\%; он выбирается на случайной половине панелей и проверяется на другой половине, по 200 разбиениям. Это даёт $k = 1.12$ на MMLU (диапазон 5--95\% по разбиениям 1.04--1.30; медианное покрытие на отложенной половине 0.951) и $k = 1.17$ на MMLU-Pro (1.05--1.30; 0.949). Калибровка проверена только для $\dE$; для стоимостей суверенитета $C$ и $C'$, сравнивающих два портфеля, тот же коэффициент применяется приближённо, поэтому их номинальный уровень 95\% не подтверждён. Все эти интервалы отражают только выборку вопросов; неопределённость выборки моделей оценивается отдельно методом складного ножа с исключением одного разработчика (\ref{sec:supp}).

\section{Результаты}

\subsection{Общие ошибки на уровне разработчика (H1)}
По заранее заданному решающему правилу H1 подтверждается на обоих бенчмарках: уверенно на MMLU и слабо на MMLU-Pro (табл.~\ref{tab:hyp}). На MMLU две модели одного разработчика дают одинаковый неверный ответ на 0.46 п.~п. большей доле вопросов (примерно 65 из 14,038), чем подобранные по точности и промпту пары разных разработчиков ($T = 0.00462$, $p \le 0.005$, где 0.005~--- наименьшее значение, достижимое при 200 нулевых панелях; 90\% интервал по методу складного ножа $[0.00162, 0.00762]$), что составляет примерно 13 стандартных отклонений нулевого распределения от среднего по N2 в заранее заданном запуске и в 1.65 раза превышает минимальный обнаружимый эффект. Результат сохраняется во всех четырёх панелях для анализа чувствительности: строгий промпт ($T=0.00991$), без крупнейшего разработчика (0.00422), без моделей с точностью ниже 40\% (0.00458) и только открытые модели (0.00595). На MMLU-Pro $T=0.00414$ ($p \le 0.005$) превышает свой минимальный обнаружимый эффект 0.00358, но подтверждение слабее: интервал по методу складного ножа включает ноль в основной панели и в трёх из пяти панелей для анализа чувствительности, панель без моделей с точностью ниже 40\% не даёт определённого результата ($T=0.00304$), а в диагностике post hoc без пары моделей одного разработчика (Writer) статистика снижается до $T=0.00151$, т.~е. ниже минимального обнаружимого эффекта 0.00358.

Эффект разработчика сохраняется и при мерах согласия из \citet{kim2025correlated} и CAPA \citep{goel2025great} после выравнивания по точности, а также без вопросов MMLU, отмеченных как ошибочные в MMLU-Redux (\ref{sec:supp}).

\begin{table}[!htbp]
\centering
\caption{Заранее заданные гипотезы: решения на обоих бенчмарках.}
\label{tab:hyp}
\small
\begin{tabular}{@{}p{3.4cm}p{4.6cm}p{4.6cm}@{}}
\toprule
Гипотеза & MMLU (70 моделей) & MMLU-Pro (58 моделей) \\
\midrule
H1 общие ошибки на уровне разработчика & подтверждена; $T=0.0046$, $p\le0.005$; устойчива в 4 панелях для анализа чувствительности & подтверждена; $T=0.0041$, $p\le0.005$; интервал по методу складного ножа включает 0; одна панель не даёт определённого результата; без Writer ниже минимального обнаружимого эффекта \\
H2 компромисс «точность--устойчивость» между правилами & опровергнута; правила голосования в пределах 0.5\pp & не определено; одна передовая модель лучшая, но неустойчивая, ниже порога \\
H3 устойчивый отбор лучше ограничения & опровергнута; выигрыш $-0.16$\pp\ $[-0.63, 0.20]$ & опровергнута; выигрыш $-2.00$\pp\ (приближённый интервал $[-4.61, -0.30]$) \\
\bottomrule
\end{tabular}
\end{table}

\subsection{Правила агрегации и отбор (H2, H3)}
На MMLU четыре правила голосования (относительное большинство, взвешенное голосование, Dawid--Skene, стекинг) различаются не более чем на 0.5 п.~п. при полной доступности и в наихудшем случае при любом размере портфеля, а лучшая отдельная модель хуже более чем на один п.~п.; H2 опровергнута. На MMLU-Pro одна передовая модель (Gemini 3 Pro preview, ошибка 9.80\%) является наиболее точным правилом при $n=3$ и 5, но при отключении её разработчика ошибка возрастает до 14.00\%, тогда как правила голосования сохраняют 12.30--13.90\%. Направление совпадает с предсказанием H2, но разрывы при полной доступности (0.5 и 0.9 п.~п.) не достигают порога в 1 п.~п.: результат не определён.

H3 опровергнута на обоих бенчмарках (табл.~\ref{tab:h3}): ни один из трёх проверенных методов отбора (жадный по точности, устойчивый к сценариям при $\lambda = 0.5$, по взаимной информации) не превзошёл ограничение «одна модель от разработчика». Это относится только к этим методам и параметрам и не доказывает, что никакой метод оптимизации устойчивости, размер ансамбля или набор сценариев не может дать лучшего результата. На MMLU ограничение имеет наименьшую ошибку при полной доступности (10.91\%), а его ошибка в наихудшем случае на отложенных данных (11.65\%) равна ошибке жадного отбора по точности (11.65\%) и немного ниже, чем у устойчивого к сценариям отбора (11.81\%). Выигрыш ограничения появляется только на MMLU-Pro, где его ошибка в наихудшем случае на отложенных данных составляет 11.10\% против 13.10\% у устойчивого к сценариям отбора (приближённые интервалы). Оценка на отложенных данных охватывает отключения отдельных разработчиков; сценарий «только API» всегда входит в обучающее семейство сценариев. Подстановочная эвристика взаимной информации хуже всех на обоих бенчмарках; на MMLU-Pro она терпит крах (57.2\%), как только отключается её сильнейший разработчик, поскольку она отбирает слабые модели, ответы которых выглядят информативными. Это касается нашей эвристики, а не оценщика из \citet{turkmen2026}. Опровержения H2 и H3 на MMLU сохраняются во всех трёх панелях для анализа чувствительности.

\begin{table}[!htbp]
\centering
\caption{Отбор портфеля ($n=5$, относительное большинство): ошибка в \%, полная доступность и наихудший случай на отложенных данных при отключении с исключением одного разработчика (leave-one-developer-out). Значения для MMLU-Pro имеют приближённые интервалы.}
\label{tab:h3}
\small
\resizebox{\textwidth}{!}{%
\begin{tabular}{@{}lrrrr@{}}
\toprule
& \multicolumn{2}{c}{MMLU} & \multicolumn{2}{c}{MMLU-Pro} \\
\cmidrule(lr){2-3}\cmidrule(lr){4-5}
Отбор & Полная & Наихудший случай & Полная & Наихудший случай \\
\midrule
Ограничение ($\le$1 от разработчика) & 10.91 & 11.65 & 9.80 & 11.10 \\
Жадный по точности & 11.06 & 11.65 & 10.90 & 14.80 \\
Устойчивый к сценариям ($\lambda=0.5$) & 11.12 & 11.81 & 10.70 & 13.10 \\
Подстановочная взаимная информация & 11.75 & 12.65 & 13.60 & 57.20 \\
\bottomrule
\end{tabular}}
\end{table}

\subsection{Стоимость отключения}
Основной портфель на MMLU (Claude 3.5 Sonnet, Yi 34B, Gemini 1.5 Pro, DeepSeek-V3, GPT-4o) ошибается на 11.06\% [8.89, 14.65] вопросов. Отключение всех моделей, доступных только через API, повышает его ошибку до 13.24\% ($\dE = 2.19$ [1.13, 5.64], RAW 1.20 [1.11, 1.60]); оставшееся голосование через правило разрешения ничьих сводится к лучшей отдельной открытой модели, DeepSeek-V3. Отключение отдельных разработчиков обходится дешевле: Anthropic 0.59 [0.22, 1.35], DeepSeek 0.46 [$-0.03$, 0.87], Google 0.38 [0.02, 0.99], OpenAI 0.27 [$-0.14$, 0.74]; для DeepSeek, чьи веса опубликованы, это значение описывает вынужденный отказ, а не утрату доступа. В терминах энергосистем наибольшая потеря при отключении одного разработчика (N$-$1) составляет 0.59 п.~п. Отключения взаимодействуют: для каждой пары из Anthropic, DeepSeek, Google и OpenAI потеря обоих обходится дороже, чем сумма двух потерь (JI от 0.60 [0.21, 1.00] до 1.26 [0.57, 1.91]; точечные оценки JI могут быть слегка смещены вверх), что согласуется с голосованием, которое деградирует быстрее по мере сокращения.

Сопоставимая стоимость суверенитета $C$ составляет 2.86 п.~п. [1.78, 3.47; калибровка приближённая, как и для всех интервалов $C$ и $C'$] на MMLU (ограничение 10.91\% против голосования пяти открытых моделей 13.77\%). Однако лучший открытый вариант на MMLU~--- одна открытая модель, DeepSeek-V3 (13.24\%), которая лучше голосования пяти открытых моделей; относительно неё стоимость составляет $C' = 2.33$ п.~п. [1.36, 3.03; приближённо]. Таким образом, ограничение системы открытыми моделями обходится на MMLU в 2.3--2.9 п.~п. На MMLU-Pro обе формы дают 5.20 п.~п. [2.51, 8.36; приближённо] (9.80\% против 15.00\%), поскольку там лучший открытый вариант~--- ансамбль открытых моделей (15.00\%), превосходящий лучшую отдельную открытую модель (16.30\%). В анализе post hoc методом складного ножа с исключением одного разработчика $C$ остаётся положительной, какой бы разработчик ни был удалён; эвристический диапазон, добавляющий эту неопределённость выборки моделей к интервалу по выборке вопросов, составляет [0.22, 5.35] на MMLU и [0.10, 10.56] на MMLU-Pro. Это не доверительный интервал с подтверждённым уровнем: он лишь показывает, что знак $C$ сохраняется при обоих источниках изменчивости (\ref{sec:supp}). Запас автономности держится на одной сильной открытой модели: без моделей DeepSeek в панели отключение моделей, доступных только через API, стоило бы 7.62 п.~п. Рис.~\ref{fig:chain} обобщает ошибку на MMLU при ограничениях доступности. На 936 вопросах по инженерным наукам, заранее заданных в качестве приближения технических вопросов энергетической отрасли, основной портфель ошибается чаще (15.81\%), а отключение моделей, доступных только через API, меняет его ошибку на 0.43 п.~п. (до 16.24\%, примерно четыре вопроса; 95\% интервал [$-1.01$, 5.87], с приближённой калибровкой [$-1.17$, 6.50]): интервал включает ноль и достаточно широк, чтобы вместить 2.19 п.~п. на всём MMLU, поэтому подвыборка не позволяет судить даже о знаке эффекта; невзвешенное голосование всех моделей теряет примерно 2 п.~п. (с 24.89\% до 26.92\%); открытый портфель с ограничением по вычислительным ресурсам из моделей не более чем со 100 млрд параметров ошибается на 19.03\% всех вопросов.

\begin{figure}[!htbp]
\centering
\begin{tikzpicture}
\begin{axis}[
  xbar, width=0.55\textwidth, height=6.2cm, xmin=0, xmax=30,
  xlabel={Ошибка (\%)}, ytick={1,2,3,4,5,6},
  yticklabels={{Российские на собственных базах, локально (4)}, {Открытые российские до 20 млрд, локально (5)}, Только открытые (голосование 5 моделей), Только открытые (лучшая), Все 70 моделей (невзвешенное голосование), Все 70 моделей (одна модель от разработчика)},
  y tick label style={font=\footnotesize}, nodes near coords, nodes near coords style={font=\small},
  bar width=9pt, enlarge y limits=0.12, axis x line*=bottom, axis y line*=left]
\addplot+[xbar, fill=gray!45, draw=black] coordinates {
  (10.91,6) (17.04,5) (13.24,4) (13.77,3) (24.85,2) (26.16,1)};
\end{axis}
\end{tikzpicture}
\caption{Ошибка на HELM MMLU в зависимости от доступного набора моделей. «Только открытые (лучшая)»~--- основной портфель после отключения моделей, доступных только через API, который сводится к DeepSeek-V3. Портфели российских моделей используют 12 локальных запусков; «Российские модели (5)» включают две дообученные модели на зарубежных базовых моделях, а вариант «на собственных базах» исключает их (post hoc). Невзвешенное голосование всех 70 моделей хуже, чем хорошо подобранный портфель.}
\label{fig:chain}
\end{figure}

\paragraph{Оценка, основанная на независимости, завышает стоимость} Повторный запуск основного конвейера на 100 панелях, смоделированных при каждой нулевой модели, показывает, как ошиблась бы оценка, игнорирующая зависимость (MMLU, описательно; табл.~\ref{tab:null}). При независимости (N1) голосование пяти моделей ошибалось бы примерно на 1\% вопросов, а отключение моделей, доступных только через API, обходилось бы в 10.2 п.~п., примерно в пять раз больше наблюдаемых 2.19; ни одна панель N1 не опустилась до наблюдаемого значения. N2 близка по стоимости отключения (2.59 [2.00, 2.91]; 12\% панелей на уровне наблюдаемого значения или ниже), но во всех 100 панелях предсказывает более точный ансамбль, чем наблюдаемый (10.00\% против 11.06\%). Реальные ошибки более зависимы, чем следует из сложности вопросов, поэтому участники голосования вносят меньший вклад, чем предполагает модель, учитывающая только сложность.

\begin{table}[!htbp]
\centering
\caption{Наблюдаемая ошибка портфеля и стоимость отключения моделей, доступных только через API, в сравнении с прогнозами нулевых моделей (MMLU; по 100 смоделированных панелей; среднее [5-й, 95-й процентиль]; описательно).}
\label{tab:null}
\small
\resizebox{\textwidth}{!}{%
\begin{tabular}{@{}lrr@{}}
\toprule
& Ошибка портфеля (\%) & $\dE$, отключены модели только через API (п.~п.) \\
\midrule
N1, независимость & 1.22 [1.05, 1.38] & 10.23 [4.42, 17.37] \\
N2, сложность вопроса и популярность неверных вариантов & 10.00 [9.64, 10.37] & 2.59 [2.00, 2.91] \\
Наблюдаемое & 11.06 & 2.19 \\
\bottomrule
\end{tabular}}
\end{table}

\paragraph{Допуск новых моделей} Модельная диаграмма допуска, подогнанная по 42 открытым моделям, оценивает шестую модель, добавляемую в открытый портфель (ошибка 13.77\%), на сетке точностей $\{0.50, 0.70, 0.85\}$ и долей копирования $\{0, 0.3, 0.6\}$. Выигрыш даёт только точность 0.85, примерно равная точности сильнейших участников: 0.96 п.~п. при ошибках, независимых от панели при заданном вопросе, и 0.40 п.~п., когда кандидат копирует ответ панели в 60\% случаев; все кандидаты с точностью 0.50 и 0.70 ухудшают голосование. Поскольку реальные ошибки более зависимы, чем допускает генератор диаграммы, её столбец с нулевым копированием, вероятно, оптимистичен.

\subsection{Отечественные модели}
Лучшая российская модель, YandexGPT-5-Lite-8B, правильно отвечает на 73.3\% вопросов MMLU; лучшая открытая модель в панели HELM, DeepSeek-V3,~--- на 86.8\% (табл.~\ref{tab:local}). Для зарубежных опорных моделей, также оценённых в HELM, локальные ответы и ответы размещённых у поставщика версий совпадают на 91\% (Qwen2.5-7B) и 87\% (Mistral-Nemo) вопросов при различии точности в пределах 0.4 п.~п.; Llama-3.1-8B совпадает на 75\% и локально на 8.7 п.~п. менее точна, из которых не более 3.8 п.~п. приходится на ответы в виде предложения, а не буквы. Несвязанные локальные пары не более похожи друг на друга, чем пары «локальный запуск--HELM» сходной точности (среднее превышение 0.0096 против 0.0138), поэтому локальное развёртывание не делает модели похожими. Квантование меняет отдельные ответы, но не точность: 8-битный прогон Qwen2.5-7B так же точен, как 4-битный (71.6\% против 71.3\%; облачный 71.4\%), но совпадает с облачным прогоном на 97.4\% вопросов вместо 91.1\% (post hoc). Три локальных запуска были повторены после того, как стали известны их первые значения точности, поскольку поставляемые с ними шаблоны чата отсутствовали или были повреждены: saiga\_yandexgpt (с 0.290 до 0.663 после восстановления официального шаблона), YandexGPT-5-Lite (с 0.729 до 0.733) и 61 слишком длинный промпт GigaChat3.1, повторно запущенный с большим контекстом. Поставляемый шаблон Llama является официальным, поэтому её сбои формата оставлены как поведение модели.

\begin{table}[!htbp]
\centering
\caption{Локальные запуски на HELM MMLU (14,038 вопросов).}
\label{tab:local}
\small
\resizebox{\textwidth}{!}{%
\begin{tabular}{@{}llrr@{}}
\toprule
Модель & Разработчик (декларированная базовая модель) & Точность & Недопустимые \\
\midrule
YandexGPT-5-Lite-8B & Yandex & 0.733 & 0.0001 \\
T-lite-it-2.1 & T-Bank (Qwen3-8B) & 0.724 & 0.0001 \\
GigaChat3.1-10B-A1.8B & Sber & 0.718 & 0.0002 \\
saiga\_yandexgpt\_8b & I.~Gusev (YandexGPT-5-Lite, по названию) & 0.663 & 0 \\
Vikhr-Nemo-12B & Vikhr (Mistral-Nemo) & 0.654 & 0.0019 \\
saiga\_nemo\_12b & I.~Gusev (Mistral-Nemo, по названию) & 0.639 & 0.0001 \\
GigaChat-20B-A3B & Sber & 0.616 & 0.0001 \\
Cotype-Nano & MTS AI & 0.539 & 0.0003 \\
\midrule
Qwen2.5-7B & Alibaba & 0.713 & 0.0027 \\
Qwen3-8B (без рассуждений) & Alibaba & 0.673 & 0 \\
Mistral-Nemo-12B & Mistral AI & 0.642 & 0 \\
Llama-3.1-8B & Meta & 0.456 & 0.0380 \\
\bottomrule
\end{tabular}}
\end{table}

Добавление любой российской модели в открытый портфель повышает его ошибку на величину от 0.08 до 0.31 п.~п., а жадный повторный отбор с российскими моделями в пуле не выбирает ни одну из них. Портфель из открытых российских моделей до 20 млрд параметров, запускаемых локально, ошибается на 24.85\% [19.69, 32.00], что на 11.1 п.~п. выше, чем у открытого портфеля. В сравнениях post hoc с подбором по размеру три российские модели 8--12B на российских базовых моделях (YandexGPT-5-Lite, saiga\_yandexgpt, GigaChat3.1) на 2.84 п.~п. [1.95, 3.74] точнее зарубежных опорных моделей 7--12B без Llama, и разница остаётся положительной при исключении любой одной модели (диапазон от 1.15 до 4.91); при всех шести российских моделях 8--12B разница составляет 1.24 п.~п. [0.56, 1.97], но исключение одной модели может изменить её знак. Добавление модели любого происхождения меняет открытый портфель на одну и ту же величину (разница $-0.01$ [$-0.08$, 0.10]). Интервалы получены повторной выборкой предметов, а не моделей. Три из четырёх зарубежных опорных моделей выпущены в 2024~г., а три российские модели на собственных базах, участвующие в сравнении, новее, так что происхождение и поколение смешаны; Qwen3-8B запускалась с отключёнными рассуждениями. Таким образом, при 8--12 млрд параметров отставания российских моделей по точности не обнаружено; это не доказывает равенства, поскольку выборка моделей мала, даты их выпуска различаются, а интервалы учитывают разброс по предметам, но не по моделям; а разрыв портфеля небольших локальных отечественных моделей согласуется с масштабом моделей, которые удалось запустить. Более крупная модель согласуется с этим, но этого не доказывает (post hoc): на всех 14\,038 вопросах MMLU T-pro-it-2.1 (32 млрд параметров, дообучение Qwen3-32B российским разработчиком) отвечает правильно на 82.0\%~--- третий результат среди 15 моделей (14 открытых моделей HELM размером 27--72 млрд и T-pro); её превосходят лишь Qwen2.5-72B (83.3\%) и Qwen2-72B (82.7\%; среднее по 14 моделям 75.2\%). Первый прогон на стратифицированной подвыборке из 2\,850 вопросов (по 50 на дисциплину) дал 84.1\%, второй результат из 15. Однако по нашему собственному правилу T-pro относится к семейству Qwen и является единственной точкой данных выше 20 млрд параметров; российские модели собственных семейств в наших запусках достигают не более 73.3\% (YandexGPT-5-Lite-8B), а GigaChat-20B-A3B (судя по названию, около 3 млрд активных параметров)~--- 61.6\%. Более крупные открытые российские модели существуют~--- GigaChat3.1-702B-A36B и GigaChat3.5-432B-A28B \citep{aisage2026gigachat}, но для локального запуска им нужны сотни гигабайт памяти, и они не проверялись; их запуск напрямую проверил бы гипотезу масштаба. Поэтому то, что разрыв объясняется масштабом, а не происхождением,~--- гипотеза, согласующаяся с данными, а не доказанный вывод. Кроме того, T-pro новее большинства моделей HELM, с которыми она сравнивается, так что масштаб и поколение смешаны. Её добавление в открытый портфель шестой моделью меняет ошибку на $-0.26$ п.~п. [$-0.56$, 0.26], а жадный повторный отбор с T-pro в пуле её не выбирает, как и предсказывает диаграмма допуска для кандидата с точностью ниже примерно 85\%. На русскоязычном тесте (ruMMLU из бенчмарка MERA \citep{fenogenova2024mera}, публичная тестовая часть; post hoc) мы использовали два промпта~--- без примеров и с пятью решёнными примерами на каждую предметную область~--- и оценивали оба на одних и тех же 9\,748 вопросах (без вопросов, взятых в качестве примеров). Голосование относительным большинством трёх российских моделей с перекрёстной подгонкой ошибается на 3.6 п.~п. реже аналогичного голосования трёх зарубежных моделей без примеров и на 5.4 п.~п. реже с решёнными примерами ($-5.42$ [$-6.93$, $-3.48$]); этот результат устойчив. Только по допустимым ответам три модели на российских базах точнее трёх зарубежных опорных моделей при обоих промптах (+4.71 п.~п. без примеров, +2.12 [+1.14, +3.50] с примерами). Для всех шести российских моделей 8--12 млрд параметров преимущество по допустимым ответам зависит от промпта (+2.15 [+1.29, +2.97] без примеров, $-0.53$ [$-1.43$, +0.62] с примерами), поэтому мы его не утверждаем. Если нарушения формата считать ошибками, шесть российских моделей 8--12 млрд уступают зарубежным опорным на 2.07 п.~п. [1.00, 3.08] без примеров и на 5.25 [2.67, 7.72] с примерами (1.49 и 2.87 при варианте парсера); именно это получила бы СППР, принимающая сгенерированные ответы. Тройка для голосования выбирается из шести российских и трёх зарубежных моделей. Решённые примеры не устраняют сбои формата ответа \citep{sclar2024format}: у некоторых моделей доля недопустимых ответов растёт, поскольку модель заново отвечает на примеры (Vikhr-Nemo с 0.22 до 0.38, Llama-3.1-8B с 0.08 до 0.45, Qwen3-8B с 0.09 до 0.20), а у GigaChat3.1 она остаётся высокой (с 0.40 до 0.36). Вариант разбора ответов post hoc, который принимает также латинскую букву с последующим текстом варианта, снижает долю недопустимых ответов T-lite с примерами с 0.157 до 0.032 и не меняет выводов. Если отключаются также зарубежные базовые семейства, российские модели на российских базовых моделях (YandexGPT, saiga\_yandexgpt и обе модели GigaChat) ошибаются на 26.16\%, и голосование почти не улучшает результат лучшей отдельной модели, поскольку остаются только два независимых разработчика.

\subsection{Наследование ошибок}
\paragraph{Уровень вопросов, локальные запуски} Шесть связанных локальных пар: три российские дообученные модели с их зарубежными базовыми моделями (Vikhr-Nemo и saiga\_nemo с Mistral-Nemo, T-lite-it-2.1 с Qwen3-8B), одна пара дообученных моделей на одной и той же зарубежной базовой модели (Vikhr-Nemo, saiga\_nemo), одна российская дообученная модель с её российской базовой моделью (saiga\_yandexgpt, YandexGPT-5-Lite) и две модели одного разработчика (GigaChat3.1, GigaChat-20B). Для каждой пары мы вычислили $x_{ij}$ минус его среднее по несвязанным локальным парам, подобранным по обеим точностям, со значимостью по нулевой модели N2 (200 панелей). Все шесть пар превышают свои пары сравнения ($p\le0.005$, нижняя граница теста; рис.~\ref{fig:inherit}a). У Vikhr-Nemo сырое превышение $x$ с её базовой моделью составляет 0.086 против 0.14--0.16 для одной и той же модели, запущенной в двух вариантах развёртывания (локально и в HELM), поэтому она сохраняет примерно половину согласия идентичной модели. Две дообученные модели на основе Mistral-Nemo имеют общее превышение 0.063 друг с другом. T-lite-it-2.1, обученная дальше и на пять п.~п. более точная, чем её базовая модель, сохраняет меньше от Qwen3-8B (0.018).

\paragraph{Масштаб популяции, Open LLM Leaderboard} Этот анализ был задан во внутреннем протоколе до вычисления каких-либо результатов, после того как стали известны числа пар и опубликованные коэффициенты из \citet{kim2025correlated}. Пары моделей рейтинга относятся к классам «родитель--потомок» (декларированная базовая модель одной модели является другой моделью; 32 пары, 8 базовых семейств), «сиблинги» (одна и та же декларированная базовая модель; 737 пар, 15 семейств), «один класс архитектуры с разной декларированной базовой моделью» (18,132 пары) и «несвязанные» (41,130 пар). Статистика~--- $P(\text{same answer} \mid \text{both wrong})$ минус её среднее по несвязанным парам, подобранным по обеим точностям в пределах $\pm0.02$, с интервалами кластерного бутстрэпа по базовым семействам. Пары «родитель--потомок» превышают подобранные несвязанные пары на 0.167 [0.114, 0.192], а сиблинги~--- на 0.101 [0.061, 0.288], во всех анализах чувствительности. Заранее заданный порядок «родитель--потомок $>$ сиблинги $>$ одна архитектура» не подтверждается; среди пар одинакового размера пары одной архитектуры (0.246) даже превышают пары «родитель--потомок» (0.155).

Проверки post hoc показывают, что это зависит от того, как размечено происхождение (рис.~\ref{fig:inherit}b; подробности~--- см. \ref{sec:supp}). При происхождении из собственной карточки каждой модели (213 из 349 моделей; 168 прямых пар «родитель--потомок» в 64 семействах) прямая связь дообучения превышает плацебо-\allowbreak сравнение с парами одной архитектуры и одного размера с разными декларированными базовыми моделями на 0.101 [0.061, 0.162] для одних только связей дообучения, тогда как производное поле базовой модели в рейтинге такого превышения не даёт.

\begin{figure}[!htbp]
\centering
\begin{tikzpicture}
\begin{axis}[name=a, title={(a) Локальные запуски, уровень вопросов:\\превышение над подобранными несвязанными парами}, title style={font=\small, align=center},
  xbar, width=0.6\textwidth, height=4.8cm, xmin=0, xmax=0.1,
  ytick={1,2,3,4,5,6}, yticklabels={Пара GigaChat, saiga\_yandexgpt--YandexGPT, T-lite--Qwen3-8B, saiga\_nemo--Mistral-Nemo, Vikhr--saiga\_nemo, Vikhr--Mistral-Nemo}, enlarge y limits=0.12,
  y tick label style={font=\scriptsize}, x tick label style={font=\scriptsize, /pgf/number format/fixed},
  bar width=6pt, axis x line*=bottom, axis y line*=left, scaled x ticks=false]
\addplot+[xbar, fill=gray!45, draw=black] coordinates {
 (0.010,1) (0.011,2) (0.018,3) (0.037,4) (0.063,5) (0.078,6)};
\end{axis}
\begin{axis}[at={(a.south west)}, anchor=north west, yshift=-1.6cm, title={(b) Рейтинг: превышение над плацебо-сравнением (одна архитектура и размер)}, title style={font=\footnotesize},
  width=0.6\textwidth, height=5cm, xmin=-0.08, xmax=0.2, ymin=0.5, ymax=4.5,
  ytick={1,2,3,4}, yticklabels={Сиблинги по карточкам,Родитель по карточке: слияние,Родитель по карточке: дообучение,Корневое поле рейтинга}, y tick label style={font=\scriptsize},
  x tick label style={font=\scriptsize, /pgf/number format/fixed}, xlabel={Превышение $P(\text{same}\mid\text{both wrong})$ над плацебо-сравнением}, xlabel style={font=\scriptsize},
  axis x line*=bottom, axis y line*=left, extra x ticks={0}, extra x tick style={grid=major}]
\addplot+[only marks, mark=*, black, error bars/.cd, x dir=both, x explicit] coordinates {
 (0.007,1) += (0.053,0) -= (0.018,0)
 (0.102,2) += (0.029,0) -= (0.030,0)
 (0.101,3) += (0.061,0) -= (0.039,0)
 (-0.020,4) += (0.045,0) -= (0.032,0)};
\end{axis}
\end{tikzpicture}
\caption{Наследование ошибок. (a) Превышение согласия в одинаковых неверных ответах $x_{ij}$ связанных локальных пар над несвязанными локальными парами, подобранными по точности (все $p\le0.005$ относительно нулевой модели N2). (b) Post hoc: превышение условного согласия связанных пар над плацебо-сравнением с парами одной архитектуры и одного размера с другой декларированной базовой моделью, с 95\% интервалами бутстрэпа по семействам; «корневое поле рейтинга» использует производное поле базовой модели в рейтинге, остальные строки~--- базовые модели, декларированные в карточках моделей. Шкалы (a) и (b) различаются.}
\label{fig:inherit}
\end{figure}

\section{Обсуждение}

\subsection{Следствия для проектирования СППР}
Результаты подкрепляют пять правил проектирования систем поддержки принятия решений на основе LLM, которые должны продолжать работать при утрате поставщиков. \emph{Во-первых, тестируйте сценарии отключения, а не только среднюю точность, и повторяйте тест при изменении портфеля.} Потери при групповом отключении измеримы стандартными мерами значимости, а положительная совместная значимость показывает, что потери нарастают по мере сокращения портфеля. \emph{Во-вторых, не оценивайте риск отключения в предположении независимости.} Оценка, основанная на независимости, завысила стоимость отключения моделей, доступных только через API, примерно в пять раз; модель, сохраняющая сложность вопросов, оказалась близка по стоимости, но завысила вклад участников ансамбля. \emph{В-третьих, ограничение «одна модель от разработчика»~--- простое и эффективное правило разнообразия.} Ни выбор среди правил голосования, ни устойчивый к сценариям отбор не улучшили его результат ни на одном из бенчмарков, а само ограничение прозрачно и дёшево; в сочетании с четвёртым правилом оно превратилось бы в ограничение на семейство моделей. \emph{В-четвёртых, считайте семейства моделей, а не бренды.} Это правило~--- гипотеза для проектирования. В проверке post hoc (\ref{sec:supp}) ограничение по семействам выбрало те же портфели, что и ограничение по разработчикам, на панели из 70 моделей, где самые точные модели и так принадлежат разным разработчикам и семействам. На 12 локальных запусках, где дообученные модели делят семейства, оно заменило saiga\_yandexgpt на Vikhr-Nemo: ошибка при полной доступности выросла на 0.49 п.~п. [0.00, 0.82], а наихудшая потеря при отключении одного семейства снизилась на 0.32 [$-0.39$, 0.25]. Таким образом, прямая проверка не показала однозначного преимущества. Точечные оценки указывают на обмен части точности на защиту от потери семейства, но ни одна из разностей статистически не установлена, поэтому это тенденция, а не устойчивый результат. Мы рекомендуем учитывать семейства при оценке разнообразия портфеля, но пока не как проверенное правило отбора. Дообученные модели сохраняют значительную часть структуры ошибок своего родителя сверх того, что разделяет их семейство моделей, а дообученные модели на одной предобученной базовой модели похожи друг на друга; портфель, считающий их независимыми поставщиками, завышает своё разнообразие. Это касается избыточности, а не доступности: веса дообученной модели остаются пригодными к использованию, когда разработчик её базовой модели прекращает обслуживание, но дообученная модель даёт лишь частичную независимую защиту от ошибок этого семейства. \emph{В-пятых, включайте новую модель, только проверив, что она улучшает портфель.} При нашем правиле агрегации и составе панели ни одна российская или зарубежная модель до 20 млрд параметров не улучшила открытый портфель, как и T-pro-it-2.1 (32 млрд) в роли шестого участника, а вклад определялся точностью, а не происхождением; в других системах, при другой маршрутизации или для отдельных типов задач менее точная модель со слабо связанными ошибками может оказаться полезной, поэтому проверять нужно и точность, и совпадение ошибок.

\subsection{Следствия для технологического суверенитета}
В рамках этой панели и процедуры отбора ограничение системы открытыми моделями обходится в 2.3--2.9 п.~п. ошибки на MMLU и примерно в 5.2 п.~п. на более сложном MMLU-Pro; это оценки для данной панели, а не общие константы. Для системы, отвечающей на 1\,000 запросов операторов сопоставимой сложности, это примерно 23--29 дополнительных неверных ответов на MMLU и 52 на MMLU-Pro при примерно 109 и 98 неверных ответах при доступности всех моделей. На MMLU лучшая открытая конфигурация~--- одна сильная открытая модель, на MMLU-Pro~--- небольшой открытый ансамбль; без моделей DeepSeek стоимость отключения на MMLU выросла бы с 2.19 до 7.62 п.~п., поэтому автономность зависит от доступа хотя бы к одному сильному открытому семейству моделей. В условиях данного исследования (вопросы с выбором ответа на английском языке, модели до 20 млрд параметров, запускаемые локально) отставания российских моделей по точности при сопоставимом размере не обнаружено, хотя происхождение и поколение здесь смешаны, и они являются ни лучшими, ни худшими участниками портфеля, чем зарубежные модели того же размера. Разрыв в 11 п.~п. у портфеля небольших локальных отечественных моделей согласуется с разрывом в масштабе моделей, которые можно запустить локально, но это остаётся гипотезой: единственная точка данных выше 20 млрд параметров, T-pro-it-2.1 (32 млрд параметров, дообучение Qwen3-32B российским разработчиком), уступает лишь двум из 14 открытых моделей HELM размером 27--72 млрд, однако по нашему правилу относится к семейству Qwen, а российские модели собственных семейств в наших запусках достигают не более 73.3\%. Если гипотеза верна, операционная автономность зависит от возможности запускать крупные открытые модели, например открытые модели GigaChat на 432 и 702 млрд параметров, и от разнообразия семейств моделей в большей степени, чем от национальной принадлежности разработчика. Диаграмма допуска количественно определяет планку для новой модели: она должна быть примерно столь же точной, как сильнейшие открытые модели, и её ошибки не должны копировать их ошибки.

\subsection{Связь с теорией надёжности}
Совпадающие отказы на трудных входных данных \citep{eckhardt1985,knight1986} соответствуют нашей нулевой модели N2, а эффекты разработчика и семейства являются зависимостями по общей причине сверх неё \citep{iaea1992ccf}; разрыв между прогнозами N2 и наблюдаемыми стоимостями отключения~--- их эффект на уровне системы. Ограничение «одна модель от разработчика»~--- простая форма принудительного разнообразия, которое в многоверсионном программном обеспечении снижает совместные отказы, не гарантируя надёжности \citep{littlewood1989}; здесь ни один из проверенных методов отбора не превзошёл ограничение. Традиционные меры разнообразия ансамблей \citep{kuncheva2003} не разделяют эти источники совпадений и не рассматривают потерю участников; меры отключения это делают.

\subsection{Угрозы валидности}
\emph{Конструктная валидность.} Точность на вопросах с выбором ответа в общих бенчмарках является приближённым показателем надёжности LLM-компонента в СППР; она не учитывает стоимость решений, контекст задачи и пользователей. Подмножество по инженерным наукам является лишь приближением отраслевых вопросов, и на нём отключение моделей, доступных только через API, изменило ошибку всего на 0.43 п.~п. (примерно четыре из 936 вопросов; 95\% интервал [$-1.01$, 5.87] включает ноль), поэтому переносимость основной оценки на технические вопросы остаётся открытым вопросом. В более общем смысле метод измеряет устойчивость точности ответов ансамбля к утрате компонентов; одинаковые значения $\dE$ могут соответствовать очень разным операционным рискам, если стоимость ошибок различается, поэтому перенос на надёжность действующей СППР требует отдельной проверки на отраслевых задачах с неравной стоимостью ошибок. Качество работы на русском языке измерено только для локальных моделей, post hoc, с промптом без примеров и с пятью примерами; при обоих сбои формата ответа смешиваются с точностью, а преимущество всех шести российских моделей по допустимым ответам зависит от промпта. \emph{Внутренняя валидность.} Примерно 6.5\% вопросов MMLU размечены ошибочно \citep{gema2025mmlu}; исключение вопросов, помеченных в этой повторной разметке, не изменило ни одного вывода, но непомеченный шум разметки может сохраняться. Понятие «разработчик» объединяет общие данные, обучение, дообучение после предобучения (post-training) и обслуживание. Модели различаются предпочтениями букв ответа, даже внутри одного семейства \citep{alzahrani2024benchmarks}; нулевая модель post hoc, в которой каждая модель выбирает неверные ответы со своими предпочтениями букв, оставляет эффект разработчика неизменным ($T=0.00466$, $p\le0.005$), так что общая склонность к буквам его не объясняет. Анализы с подбором по размеру, на собственных базовых моделях, сравнения мер и проверки устойчивости по рейтингу являются анализами post hoc, причём в сравнении с подбором по размеру российские модели новее большинства зарубежных опорных моделей; три локальных запуска были исправлены после того, как стали известны их первые значения точности, путём восстановления официальных форматов чата; заменённые запуски сохранены. \emph{Валидность статистических выводов.} H1 на MMLU-Pro опирается на одну пару разработчиков; процентильные интервалы давали недостаточное покрытие (0.91 для $\dE$ на MMLU, 0.896 на MMLU-Pro) и были откалиброваны на синтетических панелях, что предполагает, что синтетический генератор воспроизводит реальную изменчивость выборки; сравнение с нулевыми моделями описательное; анализ рейтинга использует статистики на уровне пар и самодекларируемое происхождение (базовые модели из карточек для 213 из 349 моделей), и его результаты зависят от того, как размечено происхождение; групповые сравнения локальных моделей используют повторную выборку предметов, а не моделей; метод складного ножа по разработчикам показывает, что в неопределённости мер отключения доминирует выборка моделей. \emph{Внешняя валидность.} Оценки стоимости условны относительно этой панели, момента времени и процедуры отбора: их знак устойчив внутри панели при исключении любого разработчика, но величина зависит от наличия сильной открытой модели (без DeepSeek стоимость отключения API-моделей растёт с 2.19 до 7.62 п.~п.) и не обязательно переносится на другие экосистемы моделей; повторение теста на существенно иных панелях~--- задача дальнейшей работы. Локальные запуски выполнены с квантованием (8-битная проверка показала ту же точность), крупнейшая запущенная нами российская модель имеет 32 млрд параметров, тогда как существуют открытые российские модели на 432 и 702 млрд; облачные отечественные модели не оценивались; состав панелей моделей быстро меняется.

\section{Заключение}
Стресс-тест отключения, построенный на мерах значимости теории надёжности, делает операционную автономность систем поддержки принятия решений на основе LLM измеримой. На двух бенчмарках и для этой панели моделей ограничение системы открытыми моделями обходится в 2.3--2.9 п.~п. ошибки на MMLU и в 5.2 п.~п. на MMLU-Pro; ни один из трёх проверенных методов отбора не превзошёл ограничение «одна модель от разработчика», ошибка которого в наихудшем случае на отложенных данных на MMLU равна ошибке жадного отбора по точности (11.65\%), а выигрыш появляется только на MMLU-Pro (11.10\% против 13.10\%); оценка, основанная на независимости, завышает ставки. Общие ошибки на уровне разработчика отчётливы на MMLU, но слабы на MMLU-Pro, где интервал по методу складного ножа по разработчикам включает ноль, а статистика без пары Writer (0.00151) ниже минимального обнаружимого эффекта (0.00358). Объясняется ли отставание портфеля небольших локальных отечественных моделей масштабом, а не происхождением, остаётся гипотезой, согласующейся с данными. Дообученные модели сохраняют значительную часть структуры отказов своего родителя~--- на уровне вопросов и на сотнях публичных моделей~--- сверх того, что разделяет их семейство моделей. Метод измеряет устойчивость точности ответов ансамбля к утрате компонентов; его связь с операционным риском действующей СППР ещё предстоит проверить. Следующие шаги~--- отраслевые задачи со стоимостью решений, повторение на существенно иных панелях моделей, вопросы не на английском языке, популяционные данные на уровне вопросов и запуски крупнейших открытых отечественных моделей.

\section*{Доступность данных и кода}
Все скрипты анализа, протокол с историей изменений и хешами, все сохранённые выходные данные и ответы локальных моделей (включая заменённые запуски) доступны по адресу \url{https://github.com/AndrewStrizhevskiy/withdrawal-stress-test}. Исходные результаты HELM являются публичными, а репозиторий содержит скрипты, которые их загружают и разбирают. Таблица пар Open LLM Leaderboard доступна в репозитории \citet{kim2025correlated} и не распространяется повторно.

\section*{Благодарности}
Авторы благодарят компанию Anthropic за ИИ-ассистента Claude, который помогал в подготовке кода анализа и редактировании рукописи.

\section*{Декларация об использовании генеративного ИИ и технологий с поддержкой ИИ в процессе написания}
При подготовке этой работы авторы использовали Claude (Anthropic) для помощи с кодом анализа и редактированием текста. Авторы руководили исследованием, проверили и подтвердили анализы и текст и несут полную ответственность за содержание публикации.

\section*{Декларация о конфликте интересов}
Авторы заявляют об отсутствии конфликта интересов.

\appendix
\section{Изменения протокола, анализы post hoc и заранее заданные выходные данные}
\label{app:amend}
\footnotesize
Протокол внутренний: его хеш фиксировался перед каждым этапом, но внешней регистрации с независимой отметкой времени не было. Часть решений изменена после синтетических проверок (поправки A2--A5), три локальных прогона исправлены после того, как стали известны их первые результаты точности, а несколько анализов заданы только после получения других результатов; все они перечислены ниже с датами и отмечены в тексте как post hoc. Поэтому не каждый результат следует из исходного неизменного плана.

{\renewcommand{\arraystretch}{1.25}%
\begin{longtable}{@{}>{\raggedright\arraybackslash}p{0.13\textwidth}>{\raggedright\arraybackslash}p{0.50\textwidth}>{\raggedright\arraybackslash}p{0.31\textwidth}@{}}
\toprule
Изменение & Содержание & Время относительно данных \\
\midrule
\endhead
A1 & Конкретные параметры синтетической проверки и критерии её прохождения & до проверки; без данных \\
A2 & Статистика H1 заменена попарным превышением согласия в одинаковых неверных ответах; критерий восстановления сложности заменён надёжностью повторов & после неудачи первого синтетического запуска; без реальных данных \\
A2.1 & Подбор пар сравнения по точности в H1 & в ходе реализации; без реальных данных \\
A3 & Целевая повторная проверка интервалов совместной значимости & без реальных данных \\
A4 & Промпты, различающиеся только начальной инструкцией о формате, считаются эквивалентными; промпт как переменная подбора; строгая панель как анализ чувствительности & после того, как стало известно, какие модели исключает строгое правило; до анализа какого-либо ответа \\
Точные промпты & Подбор в H1 по точному тексту промпта (задокументирован после начала запуска; реализован до него) & до просмотра какой-либо статистики \\
A5 & Повторная проверка валидации для MMLU-Pro, принятая после неудачной проверки уровня значимости (0.070): новые начальные значения генератора для уровня значимости, фиксированные сдвиги по предметам для покрытия; переход к приближённым интервалам & без реальных результатов \\
Локальные запуски & Восстановлены официальные форматы чата для трёх моделей; один повторный запуск с увеличенным контекстным окном & после первых значений точности (post hoc) \\
Локальные схемы анализа & Анализы связки (bridge), наследования и допуска для локальных запусков & заданы до локальных запусков; после результатов MMLU \\
Рейтинг & Наследование в масштабе популяции & задано до результатов; после того, как стали известны числа пар и опубликованные коэффициенты \\
Post hoc & Сравнения с подбором по размеру и на собственных базовых моделях; сравнение мер; стоимость суверенитета $C$; плацебо-сравнение, глубина и триплеты по рейтингу; происхождение по карточкам моделей; метод складного ножа по разработчикам & заданы до вычисления; после более ранних результатов \\
MMLU-Redux & Заранее заданный анализ чувствительности с исключением, выполненный поздно & задан до получения данных; выполнен после основных результатов \\
Не выполнено & Выходные данные типа S1b на MMLU-Pro & -- \\
\bottomrule
\end{longtable}}


\section{Дополнительные результаты}
\label{sec:supp}
\normalsize
В этом приложении собраны анализы устойчивости post hoc, кратко изложенные в основном тексте. Каждый из них был задан во внутреннем протоколе до того, как был вычислен.

\paragraph{Альтернативные меры согласия} Вывод не зависит от меры согласия (post hoc, табл.~\ref{tab:measures}). При том же подборе по точности и промпту условное согласие из \citet{kim2025correlated} и CAPA \citep{goel2025great} дают контрасты разработчиков, превышающие их значения по N2 на величину от 9.6 до 16.8 стандартных отклонений нулевого распределения. Без подбора по точности сырая доля одинаковых неверных ответов полностью объясняется общей сложностью и способностями (наблюдаемый контраст 0.0062 против 0.0059 при N2). Следовательно, существенным элементом контроля является обусловливание по точности; моделирование по N2 дополнительно обеспечивает калибровку каждой меры относительно одного и того же мира, в котором действует только сложность.

\begin{table}[!htbp]
\centering
\caption{Контраст разработчиков для H1 при альтернативных мерах согласия (MMLU, 70 моделей; post hoc). Наблюдаемый контраст, его среднее по отдельному набору из 200 панелей N2 и $z = $ (наблюдаемое $-$ среднее по N2)/стандартное отклонение по N2.}
\label{tab:measures}
\small
\resizebox{\textwidth}{!}{%
\begin{tabular}{@{}lrrrrrr@{}}
\toprule
& \multicolumn{3}{c}{С подбором по точности и промпту} & \multicolumn{3}{c}{Без подбора} \\
\cmidrule(lr){2-4}\cmidrule(lr){5-7}
Мера & Наблюдаемое & Среднее по N2 & $z$ & Наблюдаемое & Среднее по N2 & $z$ \\
\midrule
Превышение над N2, $x_{ij}$ (H1) & 0.00462 & 0.00005 & 14.1 & 0.00036 & 0.00013 & 0.8 \\
Сырая доля одинаковых неверных & 0.00581 & 0.00118 & 13.6 & 0.00624 & 0.00587 & 1.2 \\
$P(\text{same}\mid\text{both wrong})$ & 0.01756 & $-0.00034$ & 9.6 & 0.01422 & $-0.00178$ & 9.4 \\
CAPA & 0.01850 & $-0.00026$ & 16.8 & 0.01028 & $-0.00060$ & 10.5 \\
\bottomrule
\end{tabular}}
\end{table}

\paragraph{Исключение ошибочно размеченных вопросов} На MMLU заранее заданное исключение 374 вопросов, помеченных как ошибочные в MMLU-Redux \citep{gema2025mmlu}, не меняет результата ($T=0.00447$, $p\le0.005$, интервал по методу складного ножа $[0.00151, 0.00743]$); то же исключение меняет стоимость суверенитета $C$ с 2.86 до 2.85 п.~п., $\dE$(только API) с 2.19 до 2.26 и выигрыш H3 с $-0.16$ до $-0.12$.

\paragraph{Откалиброванные интервалы} Процентильные бутстрэп-интервалы для $\dE$ давали недостаточное покрытие на синтетических панелях (0.912 по 454 интервалам на MMLU, 0.896 по 431 на MMLU-Pro). Умножение обеих полуширин на коэффициент $k$, выбранный для покрытия 95\% на половине панелей и проверенный на другой половине (200 разбиений), даёт $k=1.12$ на MMLU и $k=1.17$ на MMLU-Pro с медианным покрытием на отложенной половине 0.951 и 0.949; по всем уровням зависимости на MMLU $k=1.14$. Калибровка расширяет, например, интервал $\dE$ (только API) с [1.24, 5.28] до [1.13, 5.64], а интервал для DeepSeek теперь включает ноль ([$-0.03$, 0.87]). Применение $k$ к $C$ и $C'$ приближённое, поскольку проверка охватывает $\dE$. Эвристический диапазон для $C$ с учётом выборки моделей складывает в квадратурах, отдельно для каждой стороны, откалиброванную полуширину по выборке вопросов и $1.96$ стандартной ошибки метода складного ножа с исключением одного разработчика (18 разработчиков на MMLU, 15 на MMLU-Pro; стандартные ошибки 1.23 и 2.21 п.~п.). Это предполагает независимость двух источников и приблизительную нормальность ошибки складного ножа; при малом числе разработчиков и одной доминирующей открытой модели ни одно из допущений не надёжно, поэтому у диапазона нет подтверждённого доверительного уровня.

\paragraph{Ограничение по семействам} Метки семейств совпадают с разработчиками, кроме Google (PaLM~2, Gemini и Gemma) и Upstage Solar Pro, отнесённой к Phi (в карточке её предварительной версии сказано, что она масштабирует Phi-3-medium). При пяти моделях ограничения по семействам, по разработчикам и совместное выбирают одинаковые портфели в обеих половинах панели из 70 моделей (ошибка 10.91\%, наихудшее отключение разработчика и семейства 11.65\%), в том числе без объединения Solar с Phi и при объединении Gemma с Gemini. На 12 локальных запусках с ранее заданными метками происхождения ограничение по разработчикам выбирает YandexGPT-5-Lite, GigaChat3.1, T-lite, Qwen2.5 и saiga\_yandexgpt; ограничение по семействам заменяет saiga\_yandexgpt на Vikhr-Nemo. Разность «по семействам минус по разработчикам» (200 бутстрэп-выборок по предметам): ошибка при полной доступности $+0.49$ [0.00, 0.82], наихудшее отключение семейства $-0.32$ [$-0.39$, 0.25], наихудшее отключение разработчика $+0.30$ [$-0.17$, 0.53]. Этот анализ разведочный.

\paragraph{Метод складного ножа по разработчикам} Выборка моделей важнее выборки вопросов: в анализе post hoc методом складного ножа с исключением одного разработчика $C$ остаётся положительной, какой бы разработчик ни был удалён (от 2.15 до 3.48 на MMLU, от 4.10 до 6.80 на MMLU-Pro), тогда как $\dE$(только API) изменяется от 1.68 до 7.62 и максимальна без DeepSeek, когда не остаётся ни одной сильной открытой модели. Наибольший выигрыш устойчивого отбора над ограничением в том же анализе методом складного ножа составляет 0.07 п.~п., поэтому опровержение H3 не зависит от какого-либо одного разработчика.

\paragraph{Происхождение в рейтинге, плацебо-сравнение и триплеты} Проверки устойчивости post hoc показывают, что этот вывод зависит от того, как размечено происхождение (рис.~\ref{fig:inherit}b). По сравнению с плацебо-\allowbreak сравнением с парами одной архитектуры и одного размера ($\pm$10\%) с разными декларированными базовыми моделями производное поле базовой модели в рейтинге не даёт обнаружимого превышения для пар «родитель--потомок» ($-0.020$ [$-0.052$, 0.025]; 31 пара, 8 семейств). Это поле фиксирует корень, часто отсутствующий или ссылающийся на саму модель, что объединяет дальних потомков. Если же происхождение берётся из собственной карточки каждой модели (\texttt{base\_model}, доступно для 213 из 349 моделей), получается 168 прямых пар «родитель--потомок» в 64 семействах. Они превышают подобранные несвязанные пары на 0.309 [0.279, 0.341], а плацебо-сравнение с парами одной архитектуры и одного размера~--- на 0.101 [0.073, 0.134]. Только связи дообучения дают такое же превышение над плацебо-сравнением (0.101 [0.061, 0.162]; 65 пар, 17 семейств), как и слияния (0.102 [0.072, 0.131]); сиблинги по карточкам не превышают своё плацебо-сравнение (0.007 [$-0.011$, 0.060]). Следовательно, прямая связь дообучения несёт общие ошибки сверх того, что объясняет класс «архитектура и размер». Похожи ли две дообученные модели одного родителя друг на друга больше, чем на родителя, зависит от родителя: в 123 триплетах по происхождению из карточек превышение для сиблингов ниже среднего превышения «потомок--родитель» (взвешенное по семействам $-0.021$, 12 семейств), за исключением случаев, когда родителем является предобученная базовая модель (Qwen2.5-14B, Qwen2.5-72B, Qwen2-72B), чьи потомки, настроенные на следование инструкциям, похожи друг на друга больше, чем на базовую модель. Этот анализ описательный, с небольшим числом семейств каждого типа. На русскоязычном тесте все пять пар с дообучением среди локальных прогонов вновь разделяют ошибки сверх нулевой модели при промпте без примеров ($p\le0.005$) и четыре из пяти~--- с пятью решёнными примерами (пара T-lite--Qwen3-8B не воспроизводится, $-0.014$), а пара GigaChat, связанная только разработчиком,~--- нет (без примеров).

\bibliographystyle{elsarticle-harv}
\bibliography{refs}

\end{document}
